% 【基础正文】所属：第4章（由 chapters/revised/04-sources.tex \input 引入）
\minitopic{从经验到世界}知觉知识具有一种表面直接性：我们通常看见桌子便知道有桌子，不先推理“我有桌子样感觉，所以最可能有桌子”。但知觉会受错觉、期待、注意和环境影响。认识论需要同时解释其直接性与可错性。 一种\term{间接实在论}{indirect realism}认为，我们直接把握的是经验或感觉材料，再由其推知外部对象。它解释了真知觉与幻觉的共同现象，却可能制造“经验帷幕”：若只能接触内部表象，凭什么跨到外部世界？\term{直接实在论}{direct realism}主张正常知觉直接呈现世界中的对象和性质；幻觉不应决定所有知觉的结构。

\minitopic{经验是证据吗}经验至少可有三种角色。第一，它造成信念；第二，它使主体有理由；第三，在合适条件下，它本身构成与事实的认知接触。因果作用不等于规范作用：药物也能造成“有苹果”的信念，却不给理由。 教条主义主张，当主体有仿佛$p$的经验且无击败者时，便获得对$p$的初步确证，无需先证明经验可靠\parencite{pryor2000}。这避免了要求每次知觉前先完成反怀疑论证明。反对者提出“认知渗透”：若经验受偏见塑造，让经验自动确证可能使偏见自我支持。 麦克道威尔强调经验必须能够进入理由空间，而不能只是盲目的因果撞击\parencite{mcDowell1994}。但若把经验概念化得过强，婴儿、动物或缺乏相应概念者的知觉知识又会受威胁。

\minitopic{析取论与共同因素}共同因素观点认为，正常知觉和不可区分幻觉共享同一种根本心理状态。\term{析取论}{disjunctivism}则说，二者至多主观不可区分：正常情形是“看见苹果”，幻觉情形是“仅仅仿佛看见苹果”，不存在决定认识地位的中性共同状态\parencite{martin2004}。 认识析取论可进一步主张，正常知觉给主体一种反思可及而又事实性的理由。它试图结合内在主义与外部世界接触，但面对“若主体无法区分两种情形，理由为何不同”的挑战。回应是：拥有区分能力不等于在每次使用理由时必须证明自己未幻觉。

\begin{argumentbox}
幻觉论证：真实知觉可能与幻觉主观不可区分；幻觉中主体不直接接触外物；若不可区分经验具有同一根本种类，则真实知觉中也不直接接触外物。析取论拒绝第三步，而非否认不可区分性。
\end{argumentbox}

\begin{comparison}
荀子论“蔽”强调局部视角、欲望和习惯可能遮蔽判断，并以“解蔽”要求扩大和校正视野\parencite{xunzi2014}。这不是一套知觉形而上学，却能补充个体经验分析：知觉错误有时不是单次感官故障，而是注意结构和长期训练的结果。
\end{comparison}

\minitopic{错觉、幻觉与误认}错觉中通常有外部对象，但其性质显得不同，如水中直棒看似弯曲；幻觉可以没有相应对象；误认则可能准确看见一个人，却把他认作另一人。三者的失败位置不同：错觉涉及性质呈现，幻觉涉及对象存在，误认涉及概念分类或身份判断。把它们统称“感官欺骗”会阻碍精确分析。 正常知觉也具有选择性。注意让某些信息进入判断，背景知识帮助识别，同时也可能导入偏见。放射科专家看见新手忽略的结构，说明训练可改善经验或至少改善经验后的分类；种族化误认则说明长期社会分类会系统扭曲判断。知觉认识论因此连接心灵哲学、认知科学和社会认识论。

\minitopic{知觉确证是否需要概念}若经验要作为理由，主体是否必须拥有表达其内容的概念？强概念论认为，进入理由空间的内容必须可用于判断；非概念论认为，经验细节可超出主体概念能力，并仍引导学习。我们能分辨两个色片不同，却说不出颜色名称，这似乎支持非概念内容；但把差异当作某项主张的理由，又似乎需要最低限度的概念组织。 可以区分经验的信息丰富度与主体在推理中使用它的能力。经验无需完全由语言概念穷尽，但信念确证要求主体以适当方式响应经验。这避免把婴儿和动物排除，也保留理由不同于纯粹刺激的规范性。

\minitopic{案例工作坊：证人辨认}证人在昏暗街道短暂看见嫌疑人，之后从一组照片中选出一人。管理员无意中用语气暗示某张照片，证人信心随之增强。后来选中的人恰是真凶。真值不消除过程缺陷：光线、呈现方式和反馈影响辨认可靠性，事后确信也不能倒推最初经验清楚。 分析应记录每个阶段：原始观看条件、记忆保存、照片阵列、管理员行为、证人陈述变化及独立证据。法律程序的双盲呈现和即时信心记录不是外在于认识论的行政细节，而是改善证据链、减少污染的制度设计。

\minitopic{技术中介的知觉}显微镜、雷达和医学影像使“看见”依赖仪器。主体不必理解全部工程原理才能获得知识，但需要适当校准、操作训练和故障监控。仪器知觉显示直接与间接并非简单二分：科学家直接看屏幕，同时通过可靠转换接触不可裸眼观察的对象。 认识主体在此可被理解为人—仪器系统。若图像算法自动增强边缘，使用者需要知道哪些特征来自样本、哪些来自处理。透明性不要求查看全部源代码，而要求获得与判断风险相称的来源说明和验证手段。
